% Integrated in Part A of the master paper.

We use genuine continuous actions after fixing a geometric base point and an actual continuous section $s:\Gamma_K\to E_U$. Changing the section by conjugation with $a\in\pi$ conjugates every centralizer and every lifted cusp below by $a$. Every subgroup is closed, every homomorphism between Tate modules is continuous and $\widehat{\mathbb Z}$-linear, every abelianization is profinite, and $\langle z\rangle$ denotes the closed subgroup generated by $z$. Put $F=K(\mu_\infty)$ and $\Gamma_K^\circ=\operatorname{Gal}(\overline K/F)$. For an abelian variety $A$, write $T(A)=\prod_\ell T_\ell(A)$, with its usual continuous Galois action. ``Divisible element'' means an element of $\bigcap_{n\geq1}nA(L)$, or of $\bigcap_{n\geq1}(L^*)^n$ in multiplicative notation.

\subsection{The arithmetic input}
\paragraph{Statement A1 (Proved; Part A(a)).}\label{res:A1}
This is the arithmetic input for cote 149, section 3, manuscript pp.~26--27, scans 54--56.
Let $K$ be a number field and $F'/K(\mu_\infty)$ finite inside $\overline K$. For every abelian variety $A/F'$ and every torus $S/F'$,
\[
 A(F')_{\rm tors}\text{ is finite},\qquad T(A)^{\Gamma_{F'}}=0,
 \qquad\operatorname{Hom}_{\Gamma_{F'}}(T(S),T(A))=0,
 \qquad\bigcap_{n\geq1}nA(F')=0.
\]
Moreover $\bigcap_{n\geq1}(F'^*)^n=\mu_\infty$.

\emph{Source input.} Ribet's appendix to Katz--Lang, Theorem 1, has the following data: $k$ is a number field, $A/k$ is an abelian variety, and $k^{\rm cyc}=k(\mu_\infty)$. Its assertion is that $A(k^{\rm cyc})_{\rm tors}$ is finite \cite[Appendix, Theorem 1, p.~315]{Ribet1981}. We use precisely this assertion, including finiteness across all primes.

\emph{Proof.} A finite set of generators of $F'/F$ and a finite set of coefficients defining $A$ lie in a number field $k_1\subset F'$ containing $K$. Enlarging $k_1$ inside $F'$ as needed gives the equality $F'=k_1(\mu_\infty)$ and an abelian variety $A_1/k_1$ with $A=A_1\otimes_{k_1}F'$. Thus Ribet applies to $A_1$ and proves the first assertion. A nonzero element of $T_\ell(A)^{\Gamma_{F'}}$ provides compatible $F'$-rational torsion points of unbounded $\ell$-power order. This contradicts that finiteness. Taking the product over $\ell$ proves the second assertion. Choose a finite extension $F''/F'$ splitting $S$. The same argument applies to $F''$. On $\Gamma_{F''}$ the module $T(S)$ is a direct sum of copies of $\widehat{\mathbb Z}(1)$ with trivial action, because $F''$ contains all roots of unity. A Galois-equivariant homomorphism $T(S)\to T(A)$ therefore has image in $T(A)^{\Gamma_{F''}}=0$. Restriction changes neither the underlying homomorphism nor its image, proving the third assertion.

For the assertion about divisible points, take $P\in\bigcap_n nA(F')$ and enlarge $k_1$ so that $P$ is defined over $k_1$. The finite nonempty sets of $n!$-division points of $P$ in $A(F')$ have a nonempty inverse limit: any finite collection of compatibility requirements is satisfied by a sufficiently high factorial division point. Choose a compatible system and let
\[
 c:\Gamma_{k_1}\longrightarrow T(A),\qquad
 c(\sigma)=(\sigma(P_{n!})-P_{n!})_n
\]
be its continuous Kummer cocycle. Its restriction to the normal subgroup $H=\Gamma_{F'}$ is zero. For $h\in H$ and $g\in\Gamma_{k_1}$, comparing $c(hg)$ and $c(g(g^{-1}hg))$ gives $h c(g)=c(g)$. Consequently $c(g)\in T(A)^H=0$ for every $g$. All the chosen division points are defined over $k_1$, so $P\in\bigcap_n nA_1(k_1)$.

Here this last intersection can be checked without a finite-generation assumption. Spread $A_1$ to an abelian scheme at one finite place $v$ of good reduction. Properness puts $A_1(k_1)$ inside $A_1(\mathcal O_{k_{1,v}})$. Reduction modulo $\mathfrak m_v^r$ is a homomorphism to a finite group for every $r$. An element divisible by every positive integer has zero image in each of these finite groups: use multiplication by the group's order. The reductions separate points, by separatedness and completeness. Hence $P=0$. This also proves directly the needed injectivity of the Kummer map $A_1(k_1)\to H^1(\Gamma_{k_1},T(A))$. In cohomological terms, the restriction map used above has kernel $H^1(\operatorname{Gal}(F'/k_1),T(A)^H)=0$; the cocycle calculation supplies its proof with profinite coefficients.

For the multiplicative assertion, the inclusion $\mu_\infty\subset\bigcap_n(F'^*)^n$ follows by taking roots of roots of unity. Let $a$ belong to the intersection and enlarge $k_1$ so that $a\in k_1^*$. A compatible root system in $F'$ gives a Kummer cocycle which is inflated from
\[
 Q=\operatorname{Gal}(k_1(\mu_\infty)/k_1)
 \subset\widehat{\mathbb Z}^{\,*}
\]
with coefficients $T=\widehat{\mathbb Z}(1)$. This subgroup is open: its index is the finite degree of $k_1\cap\mathbb Q(\mu_\infty)$ over $\mathbb Q$. Choose $m$ such that $Q$ contains the kernel of reduction modulo $m$. There is $t\in Q$ such that $t-1=du$, where $d$ is a positive integer and $u\in\widehat{\mathbb Z}^{\,*}$. Explicitly, at primes dividing $m$ take $t_\ell=1+\ell^{v_\ell(m)}$; at odd primes not dividing $m$ take $t_\ell=2$; at $2$ not dividing $m$ take $t_2=3$. All nonunit factors of $t-1$ then have finite support and finite valuation.

For a continuous cocycle $f:Q\to T$, commutativity of $Q$ gives
\[
 (t-1)f(q)=(q-1)f(t).
\]
The right side is the coboundary of $f(t)$; therefore $d[f]=0$ in $H^1(Q,T)$. The Kummer class of $a^d$ over $k_1$ is zero. The map $k_1^*\to H^1(\Gamma_{k_1},T)$ is injective. Indeed, an element with roots of every order in $k_1$ has valuation zero at every finite prime. All its roots in $k_1$ are algebraic units whose complex absolute values are uniformly bounded. There are only finitely many algebraic integers of bounded degree with all conjugates bounded, since their monic minimal polynomials have bounded integral coefficients. Repetition of roots of unbounded order shows that the element is a root of unity. Finally use an order divisible by the exponent of the finite group $\mu(k_1)$ to show that the element is $1$. Thus $a^d=1$, as required.\hfill$\square$

\paragraph{Statement A2 (Proved; the geometric input).}\label{res:A2}
This supplies the local groups used in cote 149, section 3, manuscript pp.~26--27, scans 53--56.
Let $U$ be an affine hyperbolic curve over an algebraically closed field of characteristic zero, and let $\mathcal C$ be the profinite set of cusps of its universal cover. Every $I_c$ is isomorphic to $\widehat{\mathbb Z}(1)$, and
\[
 I_c\cap I_{c'}\ne1\Longrightarrow c=c',\qquad
 N_\pi(I_c)=I_c,\qquad C_\pi(z)=I_c\quad(1\ne z\in I_c).
\]
For a curve over a characteristic-zero field, $N_{E_U}(I_c)=D_c$, its cusp decomposition group. The same normalizer identity holds with any nontrivial closed subgroup of $I_c$ in place of $I_c$. For a connected finite geometric covering with group $H\subset\pi$, compactification $X_H$, and cusp set $S_H$, the continuous quotient map induced by $U_H\hookrightarrow X_H$ fits into
\[
 \widehat{\mathbb Z}(1)\xrightarrow{\rm diagonal}
 \bigoplus_{d\in S_H}\widehat{\mathbb Z}(1)e_d
 \longrightarrow H^{\rm ab}\longrightarrow T(J_{X_H})\longrightarrow0.\tag{A.1}
\]
The image of the $d$-summand is the image of its inertia group.

\emph{Sources and proof.} The assertions with their characteristic-zero, full-profinite, hyperbolic hypotheses are Stix 2008, Lemma 10, Proposition 11, Corollaries 12--13 and (3.1), and Stix 2012, Lemma 14, Proposition 16(1)--(3) and (3.1) \cite{Stix2008,Stix2012}. We spell out the local argument that is used repeatedly. The boundary map in the localization sequence identifies a positive loop with the corresponding summand in (A.1); its only relation is the diagonal one. Once there are at least two cusps the map from any individual summand is injective, proving the required local inertia injection. Arbitrarily many cusps can be obtained in finite covers: a cusp inertia subgroup has infinite index, since an open procyclic subgroup would be the fundamental group of a finite hyperbolic cover, whose abelianization has rank at least two. The number of points over that cusp in the finite normal cover associated to $N$ is $[\pi:NI_c]$, which is unbounded as $N$ shrinks.

Two different lifted cusps can be separated at a finite level. Refine that level until there are at least three cusps. A nonzero common inertia element has a nonzero positive integral power in the subgroup defining this level, and (A.1) would put that nonzero element in two different summand lines. Those lines intersect trivially when there are at least three summands. This proves the intersection assertion. An element normalizing $I_c$, or centralizing a nonidentity element of it, conjugates $I_c$ to an inertia group meeting $I_c$ nontrivially. It therefore fixes $c$. Its geometric stabilizer is $I_c$; its arithmetic stabilizer is $D_c$. For a nontrivial subgroup $J\subset I_c$ the same argument gives $N_{E_U}(J)\subset D_c$. Conversely the cyclotomic action of $D_c$ preserves every closed subgroup of the procyclic group $I_c$, so equality holds. These arguments also prove that every abelian subgroup meeting $I_c$ nontrivially is contained in $I_c$.\hfill$\square$

\subsection{A criterion for an inertia element}
\paragraph{Statement A3 (Proved; Part A(b)).}\label{res:A3}
This supplies the inertia recognition needed for cote 149, section 3, manuscript p.~27, scans 55--56.
With the geometric data of A2, let $z\in\pi$. For every open normal subgroup $N\triangleleft\pi$ put $H_N=N\langle z\rangle$. Suppose that the image of $z$ in $H_N^{\rm ab}$ belongs to the subgroup generated by the images of the cusp inertia groups of $U_{H_N}$. Then $z\in I_c$ for some lifted cusp $c$.

\emph{Proof.} The assertion includes $z=1$, for which any cusp works. Let $z\ne1$. The quotient $Q_N=H_N/N$ is a finite cyclic group generated by $zN$. The assumed abelianized assertion, composed with the continuous surjection $H_N^{\rm ab}\to Q_N$, says that the cusp inertia subgroups generate $Q_N$. For each prime $\ell$, the $\ell$-Sylow subgroup of a finite cyclic group generated by subgroups is contained in one of those subgroups: the largest $\ell$-adic order among the finitely many generators equals the order of the Sylow subgroup. Thus the $\ell$-primary component $z_\ell$ of $z$ fixes a cusp in the finite normal cover $U_N$. The component is defined using the product decomposition of the procyclic group $\langle z\rangle$ into its Sylow subgroups.

For all $N$ the finite sets of cusps fixed by $z_\ell$ are nonempty. For any finitely many levels choose a common normal refinement; a fixed cusp there projects to fixed cusps at all those levels. Compactness consequently gives a cusp of the universal cover fixed by $z_\ell$, so $z_\ell\in I_{c_\ell}$. Select one nontrivial component $z_{\ell_0}$. Every other component commutes with it, and A2 gives $z_\ell\in I_{c_{\ell_0}}$. Finite products of the primary components are dense in $\langle z\rangle$. Since $I_{c_{\ell_0}}$ is closed, it contains $z$.\hfill$\square$

\emph{The exact comparison with Nakamura and Stix.} In Nakamura 1994, \S2.1, the data preceding (2.1.1) are a number field $k$, a smooth geometrically connected noncomplete curve of type $(g,n)$ with $n\ge1$, $2-2g-n<0$, and a full class $\mathcal D$ of finite groups, closed under subgroups, quotients and extensions. The geometric group is the pro-$\mathcal D$ completion $G$ of the displayed surface-group presentation. His (2.1.1) says that a nonidentity $z\in G$ is an inertia element exactly when there is a closed subgroup of its arithmetic normalizer mapping openly to $\Gamma_k$ and acting on $\langle z\rangle$ through the cyclotomic character. Lemma (2.1.2) says that the centralizer of a nonidentity element in the subgroup of a free generator of a free pro-$\mathcal D$ group is that generator subgroup. Corollary (2.1.3) excludes an abelian subgroup meeting two distinct free-generator subgroups nontrivially. For $Z$ the union of all conjugates of the cusp subgroups, Lemma (2.1.4) gives
\[
 z\in Z\quad\Longleftrightarrow\quad
 z\in\overline{[H,H]H^\ell\langle Z\cap H\rangle}
 \quad\text{for all open }H\ni z\text{ and all }\ell\in|\mathcal D|.
\]
These are precise restatements of (2.1.1)--(2.1.4), pp.~94--97 \cite{Nakamura1994}.

For the full profinite group, A3 implies the right side of (2.1.4) directly: for a given $H$ choose open normal $N\subset H$ and use the inclusion $N\langle z\rangle\hookrightarrow H$; its induced map on abelianizations carries boundary summands to boundary summands. Conversely an inertia element satisfies A3 by the same inclusion of its inertia subgroup at each level. Thus A3 and Nakamura's full-profinite (2.1.4) give the same recognition of $z$, through the specified inclusion and quotient maps.

Stix 2008 Theorem 20 has a characteristic-zero field $k$ satisfying
$\operatorname{Hom}_{\Gamma_k}(T(S),T(A))=0$ for every $k$-torus $S$ and abelian variety $A$, a hyperbolic $U/k$, and a nontrivial closed $I\subset\pi$. It equates containment in inertia with procyclicity and essential cyclotomic normalization, and with procyclicity and cyclotomic normalization by the entire normalizer. The 2012 manuscript numbers its refined statement Theorem 30: fix a prime $\ell$ for which every $\Gamma_k$-map $\mathbb Z_\ell(1)\to T_\ell(A)$ is zero; for nontrivial closed pro-$\ell$ $I$, containment in inertia is equivalent to $I\simeq\mathbb Z_\ell$ with essential, or full, cyclotomic normalization. Lemma 28 proves stability of that hypothesis under finite extensions by Weil restriction. Applying this statement to a nonzero $\ell$-component of a procyclic $I$ and then applying A2 to the commuting remaining components proves the 2008 assertion whenever its hypothesis holds. For number fields and the full class of finite groups it gives Nakamura (2.1.1) with $I=\langle z\rangle$. Nakamura also treats the quotient $\pi\twoheadrightarrow\pi^{\mathcal D}$ for other full classes; that scope is stated in his theorem itself, rather than inferred by an unproved passage through this quotient.

The phrase ``number field'' in Stix 2008 Lemma 19 is left as printed; A1 supplies the necessary assertion for $F'$ directly. Stix 2012 Lemma 27 explicitly permits algebraic $k/\mathbb Q$ and has the stronger prime-specific hypothesis $[kK_\ell:K_\ell]<\infty$, with $K_\ell$ the fixed field of an $\ell$-Sylow subgroup of $\Gamma_{\mathbb Q^{\rm ab}}$. In particular it applies to finite extensions of $K(\mu_\infty)$.

\subsection{The reduction from centralizers to inertia}
\paragraph{Statement A4 (Proved; Part A(c), R2 steps (1)--(3)).}\label{res:A4}
The manuscript passages are cote 149, section 2, manuscript p.~24, scan 49, and section 3, manuscript p.~27, scans 55--56.
For a number field $K$, a section $s$ and $C=C_\pi(s(\Gamma_K^\circ))$, the centralizer is trivial when $U$ is proper. When $U$ is affine hyperbolic, every nonidentity $z\in C$ belongs to an inertia subgroup $I_c$; the subgroup $s(\Gamma_K^\circ)$ fixes $c$, and $I_c\subset C$.

\emph{Proof.} Let $H\subset\pi$ be open and contain $z\in C$. A continuous action on a topologically finitely generated profinite group has open stabilizer of any open subgroup: choose an open characteristic $N\subset H$ and use the finite action on $\pi/N$. Such characteristic subgroups exist because there are finitely many subgroups of any given finite index, so the intersection of all subgroups of index at most a fixed bound is characteristic and open. Thus some open $\Gamma_{K'}\subset\Gamma_K$ stabilizes $H$ through $s$. The subgroup $H s(\Gamma_{K'})$ of $E_U$ defines a finite etale covering of $U_{K'}$ with geometric group $H$, and $s|_{\Gamma_{K'}}$ is a section for that covering. The inclusion $H\hookrightarrow\pi$ is the induced continuous fundamental-group homomorphism, with the chosen base points.

The image of $z$ in $H^{\rm ab}$ is fixed by
$\Gamma_{K'(\mu_\infty)}=\Gamma_{K'}\cap\Gamma_K^\circ$. The map in (A.1) to the Tate module of the compactified covering is Galois-equivariant. A1 therefore makes the image of $z$ there zero. If $U$ is proper, choose an open characteristic $N$ with $z\notin N$ and set $H=N\langle z\rangle$. Then $H^{\rm ab}=T(J_{X_H})$, so $z$ has zero abelianized image. But the homomorphism $H\to H/N$ factors through $H^{\rm ab}$ and takes $z$ to a nonidentity generator, a contradiction. For affine hyperbolic $U$, exactness of (A.1), applied to all $N\langle z\rangle$, is the hypothesis of A3. Hence $z\in I_c$.

For $\sigma\in\Gamma_K^\circ$, the conjugate $s(\sigma)I_cs(\sigma)^{-1}=I_{s(\sigma)c}$ contains the same nonidentity element $z$. A2 gives $s(\sigma)c=c$. The action of $D_c$ on $I_c=\widehat{\mathbb Z}(1)$ is the cyclotomic one, trivial on this subgroup; hence $I_c\subset C$. For proper genus zero, $\pi=1$. For proper genus one, $\pi=T(J_X)$ and inner conjugation is trivial, so A1 applies directly. This covers those two proper types as well.\hfill$\square$

\paragraph{Statement A5 (Proved; Part A(d), direct route).}\label{res:A5}
This repairs the proof required by cote 149, section 3, manuscript p.~27, scans 55--56.
For an affine hyperbolic geometric curve, any closed subgroup $B\subset\pi$ all of whose nonidentity elements are inertia elements is contained in a single inertia subgroup. That subgroup is unique when $B\ne1$.

\emph{Proof.} Assume $B\ne1$ and write $\mathcal T$ for the cusps $c$ with $B\cap I_c\ne1$. It is a nonempty $B$-stable set. If $B$ is not contained in one inertia subgroup, every $J_c=B\cap I_c$ is proper. The set $\mathcal T$ cannot be one $B$-orbit. Otherwise $B$ would be the union of the conjugates of a single proper closed subgroup $J_c$. A proper closed subgroup has proper image in some finite quotient $Q$ of $B$. No finite group is the union of conjugates of a proper subgroup $J$: in its action on $Q/J$, counting pairs consisting of an element and a fixed coset gives average fixed-coset count $1$. The identity fixes $[Q:J]>1$ cosets, so some element fixes none. This contradicts the proposed union.

Choose $c,c'\in\mathcal T$ in distinct $B$-orbits. There is an open normal $N\subset\pi$ such that their images at level $N$ are in distinct $B$-orbits. Indeed, failure at every level would give by compactness one element of $B$ carrying $c$ to $c'$. Set $H=BN$. This is open, contains $B$, and has distinct images of $c,c'$ among its cusps. If it has at least three cusps, keep it. If it has exactly two, hyperbolicity forces the compactification $X_H$ to have genus at least one. Every element of $B$ maps to $1$ under the continuous compactification map $H\twoheadrightarrow\pi_1(X_H)$: every nonidentity element of $B$ lies in an inertia subgroup of $H$. Choose a degree-two finite etale cover of $X_H$, using the quotient $\pi_1(X_H)\twoheadrightarrow\mathbb Z/2\mathbb Z$ supplied by $T(J_{X_H})$. Pullback to $U_H$ replaces $H$ by an open subgroup still containing $B$ and gives four cusps. The two specified images remain distinct. This proves the required three-cusp assertion without assuming that arbitrary distinct cusps can be separated by a subgroup containing $B$.

Choose $1\ne z\in B\cap I_c$ and $1\ne z'\in B\cap I_{c'}$. In (A.1) for this $H$, let their classes be $a e_c$ and $b e_{c'}$, with $a,b\in\widehat{\mathbb Z}(1)$ nonzero. Each of $zz'$ and $z^2z'$ is a nonidentity element of $B$, because its boundary image is nonzero. It must therefore lie in an inertia group. If there are at least four cusps, neither $a e_c+b e_{c'}$ nor $2a e_c+b e_{c'}$ belongs to a single summand line modulo the diagonal: a fourth coordinate forces the diagonal coefficient to vanish. If there are exactly three cusps, membership of $n a e_c+b e_{c'}$ in a single line forces that line to be the third one and forces $na=b$. Applying this for $n=1,2$ gives $a=b$ and $2a=b$, whence $a=0$, a contradiction. This proof is valid over $\widehat{\mathbb Z}$ even though that ring has zero divisors: subtraction gives $a=0$ directly. Thus $B\subset I_c$; uniqueness follows from A2.\hfill$\square$

The program \texttt{code/A\_boundary\_checks.py} verifies the derangement count in $S_2,\ldots,S_6$ and the three-cusp calculation modulo $3,5,7,9,15,25$; its exact output is attached. Modulo $2$ both products can lie in inertia lines, which records explicitly why a single finite quotient cannot replace the integral argument.

The original R2 step (4) omits the distinction between cusps and their $C$-orbits, and does not justify the existence of its covering containing $C$. Its written proof is therefore incomplete. A5 repairs both omissions and uses three cusps with the two products $zz'$ and $z^2z'$; four cusps are unnecessary. Hyperbolicity enters through A2 and, in the two-cusp stage, through the positive genus of $X_H$.

\paragraph{Statement A6 (Proved; Part A(d), literature route; Part A(e); R2).}\label{res:A6}
This answers the formula (3) of cote 149, section 3, manuscript p.~27, scan 55, and the marginal precisions on scan 56.
For a hyperbolic curve $U/K$ over a number field and a continuous section $s:\Gamma_K\to E_U$,
\[
 C_\pi(s(\Gamma_K^\circ))\ne1
 \quad\Longleftrightarrow\quad s\text{ is cuspidal}.
\]
In this event there is a unique lifted cusp $c$ fixed by $s(\Gamma_K)$ and
$C_\pi(s(\Gamma_K^\circ))=I_c$. For a proper hyperbolic curve the centralizer is $1$.

\emph{First proof.} A4 and A5 give $C\subset I_c$ and $I_c\subset C$ whenever $C\ne1$. Since $\Gamma_K^\circ$ is normal in $\Gamma_K$, its centralizer is preserved by $s(\Gamma_K)$. Thus $s(\Gamma_K)$ normalizes $I_c$, and A2 puts it in $D_c$. This is the definition of a cuspidal section, and it proves that the base cusp is $K$-rational. Conversely a cuspidal section acts on its nontrivial inertia group through $\chi$, so $I_c\subset C$. A5 and A2 force equality and uniqueness. The proper assertion is A4.\hfill$\square$

\emph{Second proof, using the opened literature.} Stix 2008 Theorem 14 assumes injectivity of points into section classes for proper genus-at-least-two curves over every finite extension of the base field. Under that hypothesis it proves cusp uniqueness, injection of the splitting classes of a cusp decomposition extension, and disjointness from finite-point sections and from other cusps. In the 2012 manuscript these are Theorem 17(1)--(3); Remark 18(1) supplies absence of divisible abelian-variety points as a sufficient hypothesis. We verify it over $F$ with A1. Explicitly, on a proper curve $Y/F'$ two points $y_1,y_2$ with equal section classes give equal classes after the morphism $y\mapsto[y-y_2]$ to $J_Y$. Its Kummer cocycle is the difference of the two point-section cocycles, as in Stix 2008 Proposition 44. Therefore $[y_1-y_2]$ is divisible by every integer and A1 makes it zero. A principal divisor $y_1-y_2$ with distinct points would define a degree-one map to $\mathbb P^1$, so this is impossible when the genus is positive.

Base change $U_F\to U$ induces a continuous homomorphism $\pi_1(U_F)\to E_U$ identifying $\pi_1(U_F)$ with the inverse image of $\Gamma_F$. This follows by descending every finite etale cover to a finite intermediate field and taking the inverse limit of the usual finite-base-change subgroup identifications. Now restrict $s$ to $\Gamma_F$. A4 makes each nonidentity element of $C$ an element of an inertia group of a cusp fixed by this restricted section. Theorem 14(1), together with the final argument of its proof applied to every finite covering and finite extension of $F$, gives uniqueness of the lifted cusp, rather than merely uniqueness of its image on $U$. Consequently all elements of $C$ lie in the same $I_c$. The previous normality argument finishes the proof. The finite-cover step here is exactly where hyperbolicity is needed in Stix's route: one takes neighborhoods whose compactifications have genus at least two. Neither proof uses surjectivity in a section conjecture.\hfill$\square$

\paragraph{Statement A7 (Proved; germs and larger subgroups, Part A(e)).}\label{res:A7}
This addresses the subgroup clause of cote 149, section 3, manuscript p.~30, scan 61.
A6 holds for every germ of section, after replacing $K$ by its finite field of definition. For a section over $K$, let $\Omega$ be any open subgroup of $\Gamma_K^\circ$. Then its centralizer in $\pi$ equals $C_\pi(s(\Gamma_K^\circ))$. More generally, let $\Delta\subset\Gamma_K$ be closed and contain such an $\Omega$. For a noncuspidal $s$ one has $C_\pi(s(\Delta))=1$. For a cuspidal $s$ at its unique lifted cusp $c$,
\[
 C_\pi(s(\Delta))=I_c^{\chi(\Delta)}
 =\prod_{\ell:\,\chi_\ell|_\Delta=1}\mathbb Z_\ell(1)
 \ \subset\ I_c=\prod_\ell\mathbb Z_\ell(1).\tag{A.2}
\]
The factors at all other primes are zero.

\emph{Proof.} The fixed field of $\Omega$ is finite over $F$ and hence is $K_1(\mu_\infty)$ for some number field $K_1\supset K$. A6 applied over $K_1$ proves the assertion for the restricted section. Cuspidality descends to $K$: choose a finite Galois extension $K_2/K$ containing $K_1$. Its cyclotomic-kernel subgroup is normal in $\Gamma_K$, and its centralizer is the same inertia group, since it contains the nonzero inertia group obtained over $K_1$ and A6 gives uniqueness. Therefore $s(\Gamma_K)$ normalizes that inertia group. This proves both descent and equality of centralizers for $\Omega$. The containment $\Omega\subset\Delta$ gives a reversed containment of centralizers. In the cuspidal case it places $C_\pi(s(\Delta))$ inside $I_c$, where the action is exactly $\chi$. For $a_\ell\in\mathbb Z_\ell(1)$, invariance is the equation $(\chi_\ell(\delta)-1)a_\ell=0$ for every $\delta\in\Delta$. The ring $\mathbb Z_\ell$ is an integral domain. The fixed subgroup is all of the factor when the character is trivial and is zero otherwise. Taking products proves (A.2).\hfill$\square$

\paragraph{Statement A8 (Proved; exceptional types and the cyclotomic assertion).}\label{res:A8}
These are the exceptional types in cote 149, section 3, manuscript p.~26, scan 54.
The cyclotomic centralizer for a proper genus-one curve is zero whenever a section exists. For geometric types $(0,0)$ and $(0,1)$ it is zero because $\pi=1$. For a form of type $(0,2)$, after a finite extension splitting the two cusps, the geometric group is $\widehat{\mathbb Z}(1)$ and its entire group centralizes every section restricted to the cyclotomic kernel. Both cusps have that same inertia group. For every hyperbolic type, including $(0,3)$ and $(1,1)$, A6 applies with the full geometric fundamental group.

\emph{Proof.} The first three assertions follow from the explicit geometric groups and A1. For $\mathbb G_m$ the Kummer covers $t\mapsto t^n$ identify the geometric group with $\widehat{\mathbb Z}(1)$; the genuine action of any section on this abelian group is the same cyclotomic action, since an inner automorphism of an abelian group is trivial. The maps from the punctured local rings at $0$ and $\infty$ induce respectively a generator and its inverse, so both inertia images are the entire group. The last assertion is A6 with its stated hyperbolicity hypothesis.\hfill$\square$

For the real base field the exact calculation is
$\Gamma_{\mathbb R}^{\circ}=1$ and therefore
$C(s)=\pi_1(U_{\mathbb C})$ for every actual section over $\mathbb R$.
Indeed $\Gamma_{\mathbb R}=\{1,c\}$ and $c$ acts on a primitive fourth
root of unity by inversion, so $\chi(c)=-1\ne1$; its cyclotomic kernel
is trivial. The centralizer of its image is the entire geometric group.
For example this gives the whole free profinite rank-two group for
the point $2$ of $U_{0,3}(\mathbb R)$, as well as for a real tangent at
$0$. The finite character calculation modulo $4$ and the nonabelian
$S_3$ quotient are recorded in \texttt{code/root\_math\_checks.py}.
At a real place of a number field, the decomposition homomorphism
$\Gamma_{\mathbb R}\to\Gamma_K$ sends $c$ to an element with the same
cyclotomic value $-1$, so that decomposition group intersects
$\Gamma_K^\circ$ trivially. This specifies the comparison with the
global calculation A6. Over $\mathbb R$, a positive-dimensional
abelian variety has $T(A)^{\Gamma_{\mathbb R}^\circ}=T(A)\ne0$;
the arithmetic vanishing input A1 thus does not hold for this base.
